(3 points + 3 points)
\begin{teilaufgaben}
\item
The function $u(x,t)$ is defined on the strip
$\Omega = [0,1] \times [0,\infty)$.

The function $u(x,t)$ satisfies in $\Omega$ the equation
\[
u_{t}(x,t) = u_{xx}(x,t)
\]
and $u(0,t) = u(1,t) = 1$ for $t \in [0,\infty)$.

One also has the initial values $u(1/3,0) = 1$ and $u(2/3,0) = 2$ .

Determine approximate values for $u(1/3,1/10)$ and $u(2/3,1/10)$.

Use the explicit method of Richardson with $\Delta x = 1/3$ and
$\Delta t = 1/30$.

\item
The domain $\Omega$ is a  rectangle with vertices:
\[
(2,0), (2,3), (-2,3), (-2,0). 
\]

The function $u(x,y)$ is defined on $\Omega$. It satisfies in $\Omega$ 
\[
\Delta u(x,y) = 10.
\]
On the boundary of $\Omega$ one has $u(x,y) = 0$. 

Determine approximate values for $u(1,1)$ and $u(-1,1)$. 

Use finite volumes with two Voronoi cells.  
\end{teilaufgaben}

\begin{loesung}
\begin{teilaufgaben}
\item
Set $U_k^{j} = \tilde u(k/3,j/3)$.
We apply the explicit method of Richardson with $r = 3/10$ and obtain. 
\begin{equation}
\begin{pmatrix*}[r]
3/10 & 4/10 & 3/10 & 0 \\
 0 & 3/10 & 4/10 & 3/10
\end{pmatrix*}
\cdot
\begin{pmatrix}1 \\  U_1^{j} \\   U_2^{j} \\  1 \end{pmatrix}
=
\begin{pmatrix} U_1^{j+1} \\  U_2^{j+1}  \end{pmatrix}.
\tag{{\bf 1P}}
\end{equation}
Therefore we have for the first time-step:
\begin{equation}
\tilde u(1/3,1/30) = 13/10, \qquad  \tilde u(2/3, 1/30) = 14/10.
\tag{{\bf 1P}}
\end{equation}
Next we have for the second time-step:
\[
\tilde u(1/3,2/30)
=
124/100,
\qquad \tilde u(2/3, 2/30)
=
125/100.
\]
Finally we have for the last time-step:
\begin{equation}
\tilde u(1/3,1/10) = 1171/1000, \qquad \tilde u(2/3, 1/10) = 1172/1000.
\tag{{\bf 1P}}
\end{equation}
\item
Set $u_1 = \tilde u(1,1)$ and $u_{-1} = \tilde u(-1,1)$.
By the finite volume method we have
\[
3 \cdot (0 - u_1) +  (0 - u_1) + 3/2 \cdot (u_{-1} - u_1)  + 2 \cdot (0 - u_1)
=
60
\]
and
\[
3/2 \cdot (u_1 - u_{-1}) +  (0 - u_{-1}) + 3 \cdot (0 - u_{-1})  + 2 \cdot (0 - u_{-1})
=
60.
\]
This leads to a system of two linear equations:
\begin{equation}
1/2 \cdot \begin{pmatrix*}[r] -15 & 3 \\  3 & -15 \end{pmatrix*}
\cdot
\begin{pmatrix} u_1  \\   u_{-1} \end{pmatrix}
=
\begin{pmatrix*}[r] 60 \\  60 \end{pmatrix*}.
\tag{{\bf 2P}}
\end{equation}
Hence
\begin{equation}
\tilde u(1,1) = \tilde u(-1,1) = -10.
\tag{{\bf 1P}}
\end{equation}
\end{teilaufgaben}
\end{loesung}
